\documentclass[10pt, aspectratio=169]{beamer}
% Preámbulo modular para presentaciones Beamer (Modelación Estadística)
% Diseño y paleta institucional del Tecnológico de Monterrey

\documentclass[aspectratio=169,xcolor={dvipsnames,table}]{beamer}

\usepackage[utf8]{inputenc}
\usepackage[spanish,mexico]{babel}
\usepackage{amsmath,amssymb,amsfonts,mathtools}
\usepackage{graphicx}
\usepackage{listings}
\usepackage{booktabs}
\usepackage{tikz}
\usetikzlibrary{matrix,arrows,decorations.pathmorphing,shapes.geometric,positioning}

% Paleta Institucional Tecnológico de Monterrey
\definecolor{TecAzul}{HTML}{0039A6}       % Azul TEC: color primario institucional
\definecolor{TecAzulOscuro}{HTML}{1A2E51} % Azul marino oscuro
\definecolor{TecRojo}{HTML}{EC2661}       % Rojo de acento (paleta miTec)
\definecolor{TecGris}{HTML}{646464}       % Gris neutro para comentarios y subtítulos
\definecolor{TecFondoBloque}{HTML}{F4F6F9}% Fondo suave para bloques

% Configuración de Tema Beamer
\usetheme{Boadilla}
\usecolortheme[named=TecAzulOscuro]{structure}
\setbeamercolor{alerted text}{fg=TecRojo}
\setbeamercolor{palette primary}{bg=TecAzulOscuro,fg=white}
\setbeamercolor{palette secondary}{bg=TecAzul,fg=white}
\setbeamercolor{palette tertiary}{bg=TecAzulOscuro!80!black,fg=white}
\setbeamercolor{title}{fg=TecAzulOscuro,bg=TecFondoBloque}
\setbeamercolor{frametitle}{fg=TecAzulOscuro,bg=TecFondoBloque}

% Bloques personalizados elegantes
\setbeamercolor{block title}{bg=TecAzulOscuro,fg=white}
\setbeamercolor{block body}{bg=TecFondoBloque,fg=black}
\setbeamercolor{block title alerted}{bg=TecRojo,fg=white}
\setbeamercolor{block body alerted}{bg=TecRojo!8,fg=black}
\setbeamercolor{block title example}{bg=TecAzul,fg=white}
\setbeamercolor{block body example}{bg=TecAzul!8,fg=black}

% Redondear bordes y añadir sombra sutil
\setbeamertemplate{blocks}[rounded][shadow=true]
\setbeamertemplate{navigation symbols}{}

% Pie de página limpio con número de diapositiva
\setbeamertemplate{footline}{
  \leavevmode%
  \hbox{%
  \begin{beamercolorbox}[wd=.7\paperwidth,ht=2.25ex,dp=1ex,left]{author in head/foot}%
    \hspace*{2ex}\insertshortauthor~~(\insertshortinstitute) \hfill \insertshorttitle
  \end{beamercolorbox}%
  \begin{beamercolorbox}[wd=.3\paperwidth,ht=2.25ex,dp=1ex,right]{date in head/foot}%
    \insertshortdate{}\hspace*{2em}
    \insertframenumber{} / \inserttotalframenumber\hspace*{2ex}
  \end{beamercolorbox}}%
  \vskip0pt%
}

% Configuración de Listings de Python para visualización pedagógica de código
\lstset{
    language=Python,
    basicstyle=\ttfamily\footnotesize,
    keywordstyle=\color{TecAzulOscuro}\bfseries,
    stringstyle=\color{TecRojo},
    commentstyle=\color{TecGris}\itshape,
    showstringspaces=false,
    breaklines=true,
    frame=lines,
    rulecolor=\color{TecAzulOscuro!30},
    backgroundcolor=\color{TecFondoBloque},
    aboveskip=0.5em,
    belowskip=0.5em,
    literate={á}{{\'a}}1 {é}{{\'e}}1 {í}{{\'i}}1 {ó}{{\'o}}1 {ú}{{\'u}}1
             {Á}{{\'A}}1 {É}{{\'E}}1 {Í}{{\'I}}1 {Ó}{{\'O}}1 {Ú}{{\'U}}1
             {ñ}{{\~n}}1 {Ñ}{{\~N}}1 {¿}{{?`}}1 {¡}{{!`}}1
}

% Metadatos institucionales por defecto
\title[Modelación Estadística]{Modelación Estadística para Ciencia de Datos e Ingeniería}
\author[J. Castillo Colmenares]{Juliho David Castillo Colmenares}
\institute[Tec de Monterrey]{Tecnológico de Monterrey \\ Escuela de Ingeniería y Ciencias}
\date{\today}

% Comandos y macros estadísticas compartidas para las presentaciones Beamer (Metropolis)

% Conjuntos numéricos básicos
\newcommand{\R}{\mathbb{R}}
\newcommand{\N}{\mathbb{N}}
\newcommand{\Q}{\mathbb{Q}}
\newcommand{\Z}{\mathbb{Z}}
\newcommand{\set}[1]{\left\{ #1 \right\}}
\newcommand{\abs}[1]{\left|#1\right|}
\newcommand{\norm}[1]{\left\| #1 \right\|}

% Comandos estadísticos y probabilísticos del curso
\newcommand{\Var}{\operatorname{Var}}
\newcommand{\cov}{\operatorname{Cov}}
\newcommand{\corr}{\rho}
\newcommand{\comb}[2]{\begin{pmatrix} #1 \\ #2 \end{pmatrix}}
\newcommand{\s}{\sigma}
\newcommand{\particion}{\mathfrak{P}}
\newcommand{\card}[1]{\#\left( #1 \right)}
\newcommand{\seq}[3]{\set{#1}_{#2}^{#3}}
\newcommand{\E}{\mathbb{E}}
\newcommand{\Prob}{\mathbb{P}}

% Entornos pedagógicos y cajas en español académico natural
\newenvironment{discusion}{%
  \begin{alertblock}{Pregunta de Discusión}%
}{%
  \end{alertblock}%
}

\newenvironment{reflexion}{%
  \begin{alertblock}{Reflexión para el Aula}%
}{%
  \end{alertblock}%
}

\newenvironment{paradoja}{%
  \begin{alertblock}{Paradoja de Exploración}%
}{%
  \end{alertblock}%
}

\newenvironment{motivacion}{%
  \begin{exampleblock}{Motivación: Ciencia de Datos en la Práctica}%
}{%
  \end{exampleblock}%
}

\newenvironment{retoclase}{%
  \begin{block}{Reto Rápido en Equipo}%
}{%
  \end{block}%
}


\title[Supuestos del ANOVA]{Sección 08.02: Supuestos del ANOVA y Diagnóstico}
\subtitle{Homoscedasticidad, Normalidad, Independencia y Alternativas No Paramétricas}
\author[J. Castillo Colmenares]{Juliho Castillo Colmenares}
\institute{Tecnológico de Monterrey}
\date{\vspace{-1.2cm}}

\begin{document}

% Slide 1: Portada
\begin{frame}[plain]
  \vspace{-0.8cm}
  \titlepage
\end{frame}

% Slide 2: Hoja de Ruta
\begin{frame}{Hoja de Ruta --- Capítulo 08: Elementos de Diseño de Experimentos (ANOVA)}
  \footnotesize
  ¿Dónde nos encontramos en el desarrollo modular del capítulo? \pause
  \begin{itemize}\itemsep=0.08em
    \item \textbf{08.01} Análisis de Varianza de un Factor (One-Way ANOVA) y Prueba $F$ \pause
    \item \textbf{08.02} Supuestos del ANOVA, Prueba de Levene/Bartlett y Diagnóstico \emph{(Hoy --- Cierre del Capítulo)}
  \end{itemize}
  \vspace{-0.05cm}
  \begin{block}{Objetivo de la Sesión}
    Auditar rigurosamente los tres supuestos sobre los que descansa la validez de toda la maquinaria construida en la Sección 08.01 --- homoscedasticidad, normalidad e independencia de los errores --- y presentar las alternativas metodológicas (ANOVA de Welch, Kruskal-Wallis) cuando dichos supuestos fallan.
  \end{block}
\end{frame}

% Slide 3: Motivación
\begin{frame}{¿Por Qué Auditar los Supuestos?}
  \footnotesize
  \begin{motivacion}
    El estadístico $F=\text{CMTR}/\text{CME}$ solo sigue una distribución $F_{k-1,N-k}$ \emph{exacta} si el modelo paramétrico subyacente ($\epsilon_{ij}\sim N(0,\sigma^2)$ i.i.d.) es correcto. Un ANOVA calculado sobre datos que violan estos supuestos puede producir decisiones completamente erróneas, aun con aritmética perfecta.
  \end{motivacion}

  \vspace{0.15cm}
  \pause
  \begin{itemize}\itemsep=0.1cm
    \item Todo científico de datos debe auditar estos supuestos \textbf{antes} de reportar resultados de un ANOVA en producción o investigación. \pause
    \item Las pruebas de diagnóstico reutilizan, en su mayoría, la misma maquinaria de descomposición de varianza de la Sección 08.01.
  \end{itemize}
\end{frame}

% Slide 4: Los tres supuestos
\begin{frame}{Los Tres Pilares del Modelo Paramétrico}
  \footnotesize
  \begin{itemize}\itemsep=0.12cm
    \item \textbf{Homoscedasticidad:} $\sigma_1^2=\sigma_2^2=\ldots=\sigma_k^2$. El ANOVA es robusto a violaciones moderadas \emph{solo si} los tamaños muestrales son iguales. \pause
    \item \textbf{Normalidad de los errores:} $\epsilon_{ij}\sim N(0,\sigma^2)$. Se audita con gráficos Q-Q y pruebas formales sobre los residuos. \pause
    \item \textbf{Independencia:} $\text{Cov}(\epsilon_{ij},\epsilon_{lm})=0$. El supuesto más crítico: su violación invalida por completo la prueba $F$ y no se corrige con transformaciones; se garantiza \emph{a priori} mediante aleatorización.
  \end{itemize}
\end{frame}

% Slide 5: Homoscedasticidad
\begin{frame}{Verificación de Homoscedasticidad: Bartlett y Levene}
  \footnotesize
  \begin{block}{Prueba de Bartlett}
    Sensible a la normalidad. Contrasta $H_0:\sigma_1^2=\ldots=\sigma_k^2$ mediante un estadístico basado en el cociente ponderado de varianzas muestrales, $\sim\chi^2_{k-1}$.
  \end{block}
  \pause

  \vspace{0.1cm}
  \begin{alertblock}{Prueba de Levene (centrada en la media) y Brown-Forsythe (mediana)}
    Estándar de oro moderno: aplica el ANOVA de un factor completo de la Sección 08.01 sobre las desviaciones absolutas $Z_{ij}=|Y_{ij}-\bar Y_{i.}|$. Robusta ante no-normalidad.
  \end{alertblock}
\end{frame}

% Slide 6: Normalidad
\begin{frame}{Diagnóstico de Normalidad de los Residuos}
  \footnotesize
  \begin{block}{Auditoría Cualitativa}
    El gráfico de probabilidad normal (\textbf{Q-Q Plot}) de los residuos estandarizados $\hat\epsilon_{ij}=Y_{ij}-\bar Y_{i.}$ revela visualmente desviaciones sistemáticas de la normalidad.
  \end{block}
  \pause

  \vspace{0.1cm}
  \begin{alertblock}{Auditoría Cuantitativa}
    La \textbf{Prueba de Shapiro-Wilk} (o Kolmogorov-Smirnov) contrasta formalmente $H_0:$ los residuos provienen de una población normal. Es la prueba de referencia para muestras pequeñas y moderadas.
  \end{alertblock}
\end{frame}

% Slide 7: Alternativas
\begin{frame}{Cuando los Supuestos Fallan: Alternativas Metodológicas}
  \footnotesize
  \begin{itemize}\itemsep=0.1cm
    \item \textbf{Homoscedasticidad violada:} ANOVA de Welch (ajusta grados de libertad del denominador ponderando por $S_i^2/n_i$) o transformaciones Box-Cox ($\log Y$, $\sqrt{Y}$). \pause
    \item \textbf{Normalidad violada (asimetría extrema, atípicos):} \textbf{Prueba de Kruskal-Wallis}, la alternativa no paramétrica basada en rangos en lugar de medias. \pause
    \item \textbf{Independencia violada:} no existe corrección algebraica a posteriori; requiere modelos de series de tiempo o efectos mixtos longitudinales.
  \end{itemize}
\end{frame}

% Slide 8: Lab Python (Bloque 1)
\begin{frame}[fragile]{Laboratorio en Python: Bartlett y Levene}
  \scriptsize
  Auditoría de homoscedasticidad para 3 \emph{pipelines} CI/CD, verificada contra \texttt{scipy.stats.bartlett} y \texttt{scipy.stats.levene} (\texttt{08.02\_anova\_assumptions.py}):
  {\lstset{basicstyle=\fontsize{3.5pt}{4.2pt}\selectfont\ttfamily,aboveskip=0.05em,belowskip=0.05em}
  \lstinputlisting[firstline=17, lastline=38]{../../code/08_diseno_experimentos/08.02_anova_assumptions.py}}
\end{frame}

% Slide 9: Lab Python (Bloque 2)
\begin{frame}[fragile]{Laboratorio en Python: Shapiro-Wilk sobre Residuos}
  \scriptsize
  Ajuste de un ANOVA de 3 grupos y auditoría de normalidad de los residuos:
  {\lstset{basicstyle=\fontsize{3.8pt}{4.5pt}\selectfont\ttfamily,aboveskip=0.05em,belowskip=0.05em}
  \lstinputlisting[firstline=41, lastline=52]{../../code/08_diseno_experimentos/08.02_anova_assumptions.py}}
\end{frame}

% Slide 10: Lab Python (Bloque 3)
\begin{frame}[fragile]{Laboratorio en Python: Kruskal-Wallis}
  \scriptsize
  Comparación entre ANOVA paramétrico y Kruskal-Wallis sobre datos asimétricos (exponenciales):
  {\lstset{basicstyle=\fontsize{3.6pt}{4.3pt}\selectfont\ttfamily,aboveskip=0.05em,belowskip=0.05em}
  \lstinputlisting[firstline=55, lastline=70]{../../code/08_diseno_experimentos/08.02_anova_assumptions.py}}
\end{frame}

% Slide 11: Salida Terminal
\begin{frame}[fragile]{Laboratorio en Python: Salida en Terminal}
  \scriptsize
  Salida estándar generada al ejecutar el script del laboratorio en Python:
  \vspace{0.02cm}
  \begin{block}{Salida Estándar en Consola}
    \fontsize{3.8pt}{4.6pt}\selectfont
    \begin{verbatim}
=== Block 1: Bartlett and Levene ===
Bartlett: B=4.4660, p=0.1072
Levene (manual): W=2.9880, p=0.0885 (matches scipy.stats.levene center='mean')

=== Block 2: Shapiro-Wilk ===
W=0.9781, p=0.7721 -> no evidence against normality

=== Block 3: Kruskal-Wallis ===
Parametric ANOVA: F=3.5514, p=0.0376
Kruskal-Wallis:    H=11.9985, p=0.0025
    \end{verbatim}
  \end{block}
\end{frame}

% Slide 16: Cierre
\begin{frame}{Cierre de Sección y de Capítulo: Supuestos del ANOVA}
  \begin{reflexion}
    Con esta sección cerramos el Capítulo 08: construimos el ANOVA de un factor como generalización de la comparación de medias, y ahora auditamos rigurosamente los supuestos que garantizan su validez. Ningún resultado de un ANOVA es confiable sin verificar homoscedasticidad, normalidad e independencia de los errores.
  \end{reflexion}

  \vspace{0.2cm}
  \pause
  \begin{block}{Lecciones Clave}
    \begin{itemize}\itemsep=0.08cm
      \item \textbf{Levene $\equiv$ ANOVA sobre desviaciones:} ninguna maquinaria nueva, solo una transformación de variable.
      \item \textbf{Robustez condicionada:} el ANOVA tolera heteroscedasticidad moderada solo con tamaños muestrales iguales.
      \item \textbf{Alternativas:} Welch para heteroscedasticidad, Kruskal-Wallis para no-normalidad severa.
    \end{itemize}
  \end{block}
\end{frame}

% Slide 17: Perspectiva Modular
\begin{frame}{Perspectiva Modular --- El Próximo Reto}
  \scriptsize
  \begin{retoclase}
    Con el Capítulo 08 completo, dominamos la comparación de $k$ medias mediante un único factor categórico. El Capítulo 09 (Regresiones Lineales y Múltiples) generaliza esta lógica a variables predictoras \textbf{continuas}: en lugar de comparar grupos discretos, modelaremos la relación funcional entre una variable respuesta y uno o varios regresores mediante mínimos cuadrados ordinarios (OLS).
  \end{retoclase}

  \vspace{0.08cm}
  \pause
  \begin{alertblock}{Preparación Recomendada}
    \begin{itemize}\itemsep=0.06cm
      \item Repasa la interpretación geométrica de la descomposición de sumas de cuadrados $\text{SCT}=\text{SCTR}+\text{SCE}$: en regresión aparecerá como $\text{SCT}=\text{SCReg}+\text{SCE}$.
      \item Reflexiona sobre cómo el coeficiente de determinación $R^2$ del ANOVA (Sección 08.01) anticipa el $R^2$ de un modelo de regresión.
    \end{itemize}
  \end{alertblock}
\end{frame}

\end{document}
